\documentclass[10pt,a4paper]{amsart}
\usepackage[margin=19mm]{geometry}
\usepackage{amsmath,amssymb,mathtools}
\usepackage[scr=rsfs,cal=euler]{mathalfa}
\usepackage{xfrac}
\usepackage[unicode,hidelinks,bookmarks=false]{hyperref}
\newcommand{\ie}{i.e.\ }
\newcommand{\cf}{cf.\ }
\newcommand{\resp}{respectively\ }
\newcommand{\Subsection}{\subsection}
\providecommand{\nobreakdash}{\nobreak\hbox{-}}
\setlength{\emergencystretch}{2em}
\multlinegap0pt
\usepackage{amssymb}
\usepackage{mathtools}
\usepackage[shortlabels]{enumitem}
\usepackage{xcolor}
\newtheorem{theorem}{Theorem}
\newtheorem{corollary}{Corollary}
\newtheorem{lemma}{Lemma}
\newcommand{\Char}{\operatorname{Char}}
\newcommand{\D}{\xi}
\newcommand{\Horz}{\mathcal H}
\newcommand{\R}{\mathbb R}
\newcommand{\dd}{\mathrm d}
\newcommand{\g}{\mathfrak g}
\newcommand{\rank}{\operatorname{rank}}
\newcommand{\z}{\mathfrak z}
\title{Dynamical Geometry of Principal Bundle Constrained Systems:\\Compatible Pairs, Contact Degeneration, and Adapted Metrics}
\author{Dongzhe Zheng}
\date{Source: arXiv:2505.16766v3 (CC BY 4.0)}
\begin{document}
\maketitle

\noindent\emph{Source and licence.} Dynamical Geometry of Principal Bundle Constrained Systems:\\Compatible Pairs, Contact Degeneration, and Adapted Metrics. Version arXiv:2505.16766v3 (CC BY 4.0), distributed by arXiv under the Creative Commons Attribution 4.0 International licence, http://creativecommons.org/licenses/by/4.0/, as recorded in the arXiv licence metadata for this exact version and shown on the abstract page https://arxiv.org/abs/2505.16766v3. Full licence notice: \texttt{licenses/arxiv-2505.16766-CC-BY-4.0.txt}.\par\smallskip

\noindent\textit{Editorial scope.} Counted unit: the article's theorem thm:critical (source0.tex lines 393--417). The article's complete original proof of it is delivered with it (source0.tex lines 419--579, one \texttt{proof} environment of 161 lines). No mathematics is written, repaired or re-derived: the statement and the proof are the article's own lines, in the article's own notation, and no retained line is substituted or re-flowed.  Retained uncounted as the source-local context the proof needs: lines 234--257; lines 287--297; lines 347--359; lines 353--357.  Nothing was omitted from any retained span, so the package carries no omission record.  No retained line contains a citation, so the wrapper carries no bibliography and no bibliography entry.  No retained line contains a figure environment, an image-inclusion command or drawing code, so no image is delivered and the package declares no asset dependency.  Editorial formatting only, in addition: an AMS \texttt{amsart} preamble that restates the article's own declarations and macros used by the retained lines, namely the environments \texttt{theorem}, \texttt{corollary}, \texttt{lemma} on separate counters, together with the standard packages those lines need.  The retained environments print in this document as Theorem 1, Corollary 1, Lemma 1 in order of appearance, whereas the article prints the same items with the numbering of its own full document.  The complete native source of the article as distributed in the arXiv e-print for this exact version is bundled unchanged as \texttt{sources/178592/source0.tex}.
\begin{theorem}[Levi decomposition]\label{thm:rank}
Under the preceding hypotheses, $\dd\beta=0$, and
\begin{equation}\label{eq:dalpha}
 \dd\alpha=B(u,F_A)+\tfrac12B(u,[A\wedge A]).
\end{equation}
Writing $\Horz=\ker A$, there is a direct sum
\begin{equation}\label{eq:splitD}
 \D_p=\Horz_p\oplus(\mu^\perp)^\#_p.
\end{equation}
For $v,w\in T_xM$ and $a,b\in\mu^\perp$,
\begin{equation}\label{eq:leviblock}
 \dd\alpha_p(v^H+a^\#,w^H+b^\#)
 =\beta_x(v,w)+B(\mu,[a,b]).
\end{equation}
In particular,
\begin{align}
 \Char(\D)_p
 &= (\ker\beta_x)^H\oplus(\g_\mu\cap\mu^\perp)^\#_p,
 \label{eq:char}\\
 \rank(\dd\alpha|_{\D_p})
 &=\rank\beta_x+\dim\mathcal O_\mu.
 \label{eq:rank}
\end{align}
\end{theorem}

\begin{corollary}[Three geometric regimes]\label{cor:regimes}
The distribution $\D$ has the following properties.
\begin{enumerate}
\item It is Frobenius integrable if and only if $u(p)\in\z(\g)$ and
      $\beta=0$.
\item At $p$, its sections and their brackets span $T_pP$ if and only if
      $u(p)\notin\z(\g)$ or $\beta_{\pi(p)}\ne0$.
\item The form $\alpha$ is contact if and only if $\beta$ is symplectic and
      $\dim\g_{u(p)}=1$.
\end{enumerate}
\end{corollary}

\begin{lemma}[Connections preserving the adjoint field]\label{lem:scalar-source}
The map
\begin{equation}\label{eq:scalar-source}
 b\longmapsto A_b=A_0+\frac{u}{c_u}\,\pi^*b
\end{equation}
is an affine bijection from $\Omega^1(M)$ onto $\mathcal A_u$. Moreover,
\begin{equation}\label{eq:scalar-source-curvature}
 F_{A_b}=F_{A_0}+\frac{u}{c_u}\pi^*(\dd b),\qquad
 \alpha_b=\alpha_0+\pi^*b,\qquad
 \beta_b=\beta_0+\dd b,
\end{equation}
and $F_{A_b}=(u/c_u)\pi^*\beta_b$.
\end{lemma}

\begin{equation}
 F_{A_b}=F_{A_0}+\frac{u}{c_u}\pi^*(\dd b),\qquad
 \alpha_b=\alpha_0+\pi^*b,\qquad
 \beta_b=\beta_0+\dd b,
\end{equation}

\begin{theorem}[Critical cost and sharp Sobolev threshold]\label{thm:critical}
There exist constants $0<r_0<R_*/e$ and $c_1,c_2>0$, depending on $(A_0,u)$, the chart,
and norms, such that for $0<r<r_0$ and every $T>0$,
\begin{equation}\label{eq:critical-two-sided}
 \frac{c_1}{T\log(R_*/r)}\leq\mathcal C_T(r)
       \leq\frac{c_2}{T\log(R_*/r)}.
\end{equation}
The upper bound is realized up to its constant by paths with $A_t=A_0$
over $M\setminus B_{R_*}$.
The lower bound holds for every smooth path in $\mathcal A_u$ with
the stated endpoints, including paths that lose contactness before $T$.

The paths realizing the upper bound also satisfy, for a constant $C$
depending only on the same fixed data,
\begin{equation}\label{eq:critical-uniform}
 \sup_{0\leq t\leq T}
 \left(\|A_t-A_0\|_{H^{n+1}}^2+
       \|F_{A_t}-F_{A_0}\|_{H^n}^2\right)
       \leq\frac{C}{\log(R_*/r)}.
\end{equation}
Consequently, the locus of connections $A\in\mathcal A_u$ for which
$B(u,A)$ is contact is not open in the relative $H^s$ topology for any
$s\leq n+1$, including the critical exponent. It is open for every
$s>n+1$.
\end{theorem}

\begin{proof}
\smallskip\noindent
\emph{Construction.}
Choose a smooth nondecreasing function $\chi:\R\to[0,1]$ such that
$\chi=0$ on $(-\infty,0]$, $\chi=1$ on $[1,\infty)$, and $0<\chi<1$ on $(0,1)$.
For $L\geq1$ and $\rho>0$, set
\begin{equation}\label{eq:log-cutoff}
 r_L=R_*e^{-L},\qquad
 \chi_L(\rho)=\chi\!\left(\frac{\log(R_*/\rho)}{L}\right).
\end{equation}
Extend $\chi_L$ by one at the origin. It is smooth, nonincreasing, equals
one for $\rho\leq r_L$, and vanishes for $\rho\geq R_*$. Put
\[
 b_L=\chi_L\lambda_{\mathrm{std}}
\]
on the chart and extend it by zero. The flatness of $\chi$ at its endpoints
makes this a smooth global one-form. Define
\begin{equation}\label{eq:critical-path}
 A_{L,t}=A_0-\frac{t}{T}\frac{u}{c_u}\pi^*b_L,
 \qquad 0\leq t\leq T.
\end{equation}
Lemma~\ref{lem:scalar-source} places the entire path in $\mathcal A_u$.
In the chart its paired curvature is
\[
 \beta_{L,t}=\dd(\psi_{L,t}\lambda_{\mathrm{std}}),\qquad
 \psi_{L,t}=1-\frac{t}{T}\chi_L.
\]
Expansion of the exterior power gives
\begin{equation}\label{eq:radial-volume}
 \beta_{L,t}^{\,n}
 =\psi_{L,t}^{\,n-1}
   \left(\psi_{L,t}+\frac{\rho}{2}\partial_\rho \psi_{L,t}\right)
      \omega_0^n.
\end{equation}
Indeed,
$\lambda_{\mathrm{std}}=\iota_{Z_0}\omega_0$, where
$Z_0=\frac12\sum_j(x_j\partial_{x_j}+y_j\partial_{y_j})$, and
$\dd\psi(Z_0)=(\rho/2)\psi'$. For $t<T$, both factors in
\eqref{eq:radial-volume} are positive, since $\psi_{L,t}>0$ and
$\partial_\rho \psi_{L,t}\geq0$. Thus $\beta_{L,t}$ is symplectic and
Corollary~\ref{cor:regimes} gives contactness. At $t=T$ it vanishes
identically on $B_{r_L}$. The one-form $\alpha_{L,T}$ remains nowhere zero:
its value on the vertical vector $u^\#$ is still $c_u$. Its Levi rank on
that ball is exactly $\dim G-1$, by Theorem~\ref{thm:rank}.

\smallskip\noindent
\emph{The critical estimate.}
For $r_L<\rho<R_*$ and every integer $j\geq1$, differentiation of
\eqref{eq:log-cutoff} gives
\begin{equation}\label{eq:cutoff-derivatives}
 |\partial^j\chi_L|\leq\frac{C_j}{L\rho^j},\qquad
 |\chi_L|\leq\frac{C\log(R_*/\rho)}{L}.
\end{equation}
Here $\partial^j$ denotes coordinate derivatives of total order $j$.
Every differentiated term contains at least one derivative of $\chi$ and hence
one factor $L^{-1}$. Since the coefficients of $\lambda_{\mathrm{std}}$ are linear,
\begin{equation}\label{eq:primitive-derivatives}
 |\partial^k b_L|\leq\frac{C_k}{L\rho^{k-1}}
 \qquad(k\geq2)
\end{equation}
on the annulus. In dimension $2n$, the top derivative therefore obeys
\begin{align}
 \int_{r_L<\rho<R_*}|\partial^{n+1}b_L|^2\,\dd x
 &\leq\frac{C}{L^2}
       \int_{r_L}^{R_*}\rho^{2n-1-2n}\,\dd\rho \notag\\
 &=\frac{C}{L^2}\log(R_*/r_L)=\frac{C}{L}.
 \label{eq:critical-top-derivative}
\end{align}
The derivatives of orders $2$ through $n$ contribute $O(L^{-2})$.
For orders zero and one, use the second bound in
\eqref{eq:cutoff-derivatives} and the finiteness of
\[
 \int_0^{R_*}\log(R_*/\rho)^2\rho^{2n-1}\,\dd\rho
\]
and of the integral with two additional powers of $\rho$. The contribution
of the inner ball is exponentially small at these orders, and all its
derivatives of order at least two vanish because $b_L=\lambda_{\mathrm{std}}$ there.
Equivalence of coordinate and intrinsic norms on this fixed chart yields
\begin{equation}\label{eq:critical-b-norm}
 \|b_L\|_{H^{n+1}}^2+\|\dd b_L\|_{H^n}^2\leq\frac{C}{L}.
\end{equation}
Multiplication by the fixed smooth section $u/c_u$ is bounded in these norms.
Equations~\eqref{eq:critical-path} and
\eqref{eq:scalar-source-curvature} prove
\eqref{eq:critical-uniform}, and
\[
 \int_0^T\|\dot A_{L,t}\|_{H^{n+1}}^2\,\dd t
 \leq\frac{C}{TL}.
\]
Choosing $L=\log(R_*/r)$ proves the upper bound in
\eqref{eq:critical-two-sided}. Since the terminal connections are noncontact
and converge to $A_0$ in $H^{n+1}$, and hence in every lower $H^s$ norm,
this also proves the claimed failure of openness at and below the endpoint.

\smallskip\noindent
\emph{The lower bound.}
Let $A_t$ be any smooth path in $\mathcal A_u$ starting at $A_0$ and
satisfying $\beta_T|_{B_r}=0$. Write
$\delta\beta=\beta_0-\beta_T$, so $\delta\beta=\omega_0$ on $B_r$.
Choose a fixed chart cutoff $\chi_0$ equal to one on $B_{R_*}$ and supported
in the larger Darboux chart. Extend
\[
 w=\frac{\chi_0}{n}\sum_{j=1}^n
                  \delta\beta(\partial_{x_j},\partial_{y_j})
\]
by zero to $\R^{2n}$. Then $w=1$ on $B_r$, and
\begin{equation}\label{eq:scalar-trace-bound}
 \|w\|_{H^n(\R^{2n})}\leq C\|\delta\beta\|_{H^n(M)}.
\end{equation}
Take $\varphi\in C_c^\infty(B_1)$ with $\varphi\geq0$ and
$\int\varphi=1$, and set
$\varphi_r(x)=r^{-2n}\varphi(x/r)$. Then $\int w\varphi_r=1$.
The Fourier formula for the negative Sobolev norm gives, up to its fixed
normalization,
\[
 \|\varphi_r\|_{H^{-n}}^2
 =\int_{\R^{2n}}(1+|\xi|^2)^{-n}
                         |\widehat\varphi(r\xi)|^2\,\dd\xi.
\]
The region $|\xi|\leq1$ contributes a bounded amount. On
$1<|\xi|\leq r^{-1}$ the radial bound is
$C\int_1^{1/r}\rho^{-1}\,\dd\rho$. On $|\xi|>r^{-1}$, rapid decay
$|\widehat\varphi(r\xi)|\leq C_N(r|\xi|)^{-N}$ gives a bound independent
of $r$. Absorbing the fixed chart scale into the constant, for small $r$ we
obtain
\begin{equation}\label{eq:critical-dual-bump}
 \|\varphi_r\|_{H^{-n}}^2\leq C\log(R_*/r).
\end{equation}
Sobolev duality and~\eqref{eq:scalar-trace-bound} imply
\begin{equation}\label{eq:critical-curvature-lower}
 \|\delta\beta\|_{H^n(M)}^2\geq\frac{c}{\log(R_*/r)}.
\end{equation}
By Lemma~\ref{lem:scalar-source},
\[
 \pi^*\delta\beta
 =-\dd\!\left(\int_0^T B(u,\dot A_t)\,\dd t\right).
\]
The expression in parentheses is basic. Boundedness of this fixed
first-order operator and Cauchy--Schwarz in time give
\[
 \|\delta\beta\|_{H^n}^2
 \leq CT\int_0^T\|\dot A_t\|_{H^{n+1}}^2\,\dd t.
\]
Combining this with~\eqref{eq:critical-curvature-lower} proves the lower
bound. The argument applies to every path with the prescribed terminal curvature.

\smallskip\noindent
\emph{Supercritical openness.}
Let $s>n+1$ and $A'=A_0+a\in\mathcal A_u$. Its paired curvature satisfies
$\pi^*(\beta'-\beta_0)=\dd B(u,a)$. Sobolev embedding on the
$2n$-dimensional base yields
\[
 \|\beta'-\beta_0\|_{C^0}\leq C_s\|a\|_{H^s}.
\]
The minimum singular value of
$\beta_0^\flat:TM\to T^*M$ has a positive lower bound on compact $M$.
If the displayed norm is smaller than that bound, $\beta'$ is nondegenerate.
The centralizer of the unchanged field $u$ is still one-dimensional, so
Corollary~\ref{cor:regimes} makes $B(u,A')$ contact. The same argument at
each connection in this locus proves openness.
\end{proof}

\end{document}
